\subsection{
  Replies to comments received from Boris Popov - Dec. 28, 2012 }

\begin{enumerate}

%........................................................................................................
\vspace{5 mm}
\item
{\em
   ``What would be the best fit point if you allow the fit to go into the 
   unphysical region of $sin^{2}(2\theta_{23}) > 1$ ?''
}\\
\\
{\color{red}Note that the reply to this comment was modified between versions 1 and 2 of this addendum.}\\

An answer can not be provided in a 3-flavour
framework due to the existence of a $\sqrt{1-sin^{2}(2\theta_{23})}$ term in the oscillation probabilities.
The upper limit on $sin^{2}(2\theta_{23})$ was set to 1.3 and a `2+8'
fit to the Run 1+2+3 dataset was performed in a simple 2-flavour framework.
The fit yields the following best-fit values:
$sin^{2}(2\theta_{23})$ = 1.025 and
$|\Delta m^2_{32}| = 2.47 \times 10^{-3} eV^2/c^4$.
The same 2-flavour fit, in the presence of the physical boundary, yields
$sin^{2}(2\theta_{23})$ = 1.000 and
$|\Delta m^2_{32}| = 2.47 \times 10^{-3} eV^2/c^4$
as the best-fit values.

The best-fit values quoted above will be added in version 3 of T2K-TN-141.

%........................................................................................................
\vspace{5 mm}
\item
{\em
   ``How do you explain the fact that in the `2+N' configuration most of 48 systematic parameters 
   do not move by more than $\pm$0.1 $\sigma$? Could this point to some problems with the 
   penalty term definition in the Eq. (3) ?''
}\\
\\
We think that the definition of the penalty term used in this analysis is fairly standard in HEP.
It is clear though that, in the `2+N' fit, the vast majority of the systematic parameters do not
actually have any impact on the fit. This is because the single $\mu$-like ring event reconstructed 
energy  spectrum has little sensitivity to these parameters. The changes that non-nominal values of these
systematic parameters typically induce in the reconstructed energy spectrum are small, and they cannot 
offset the penalty term which disfavours such deviations. As a result, the fitted values are not distributed as 
one might expect, and, for all but the most important parameters, they are close to the nominal values.

%........................................................................................................
\vspace{5 mm}
\item
{\em
   ``Quoting the best fit point as 
   $|\Delta m^2_{32}| = 2.442 \times 10^{-3} eV^2$ do you think the 
   last digit is really significant? or we can round it off as $2.44 \times 10^{-3} eV^2$ ?''
}\\
\\
Indeed, as the error on the $|\Delta m^2_{32}|$ is, currently,
of the order of a few $\times$ 10$^{-4} eV^2/c^4$,
we do not think that the last digit is significant.
Had we been preparing a document to be shown outside the T2K collaboration,
the best-fit value would have been rounded to $2.44 \times 10^{-3} eV^2/c^4$.

However, the purpose of this technical note is somewhat different.
For the benefit of the collaboration, we strive to document all aspects of this analysis,
many of which are too detailed to be shown outside T2K.
In some situations, this forces us to use more decimal places.
For example, in Table 3, it is necessary to quote at least 5 decimal places in the numbers of events
so that the contribution of $\bar{\nu_{e}}$ CC coherent events can be seen.
Also in Table 12, we need to quote the best-fit values of $|\Delta m^2_{32}|$ down to the level of
10$^{-6} eV^2/c^4$ in order to see the difference between the `2+N' and `2+8' fits.
We prefer to keep a larger number of decimal places
so that the collaboration can gauge the size of all contributions to the spectrum,
the effect of different fitting procedures, etc.
But any results chosen to be shown outside the collaboration would have to be appropriately rounded.

%........................................................................................................
\vspace{5 mm}
\item
{\em
   ``When comparing your results with the MINOS and SK atmospheric measurements (Figs.74 and 75), 
   I think it would be important to mention explicitly in the figure whether the corresponding fits 
   are performed with 2- or 3-flavours, e.g. for the 'SK zenith' you show only 3-flavour result,
   while Itow-san in his talk at Neutrino'2012 has also presented the 'SK zenith 2-flavour' contour (p.22). 
   In other word, my suggestion would be to include in the insert the number of flavours used in the 
   fit for each result, e.g. 'T2K 3-flavour', 'MINOS 2-flavour', etc.''
}\\
\\
Since our analysis is carried out in a 3-flavour framework, we try wherever possible to compare our results
with the results of other 3-flavour analyses.
Only the SuperK zenith-angle analysis was upgraded to a 3-flavour framework.
The SuperK L/E and the MINOS combined accelerator+atmospheric analyses were done in a 2-flavour framework.
As suggested, a comment will be added in the figure legends in version 2 of this note to avoid confusion.

%........................................................................................................
\vspace{5 mm}
\item
{\em
   ``Let me also mention that contrary to want is mentioned in the TN141 the NA61/SHINE long target results 
   (ref. [18]) are not yet used for the official T2K neutrino flux tuning. 
   Please, update this for the next versions of the TN141 to avoid confusion.''
}\\
\\
In version 2 of this note, the NA61 reference will be replaced with:
N Abgrall et al. Measurements of Cross Sections and Charged Pion Spectra in Proton-Carbon Interactions 
at 31 GeV/c. Phys.Rev., C84:034604, 2011; 
N. Abgrall et al. Measurement of production properties of positively charged kaons in proton-carbon 
interactions at 31 GeV/c. arXiv:hep-ex/1112.0150, 2011.

\end{enumerate}
